\section{Introduction}
\label{sec:introduction}

Morphing transforms one 3D object into another through a sequence of intermediate states, with applications in animation, visual effects, and 3D content creation. Given two textured 3D objects, a successful transition must jointly evolve geometry and appearance, produce plausible intermediate states, and progress smoothly between them. %However, two objects do not define a unique morphing trajectory, as multiple transformations may be equally valid. 
In addition, learning such transitions is further complicated by the lack of ground-truth morphing trajectories, since existing 3D asset collections typically provide individual objects rather than intermediate sequences: for this reason, recent generative approaches use pretrained 3D models to generate intermediate states.
% For this reason, recent generative approaches to textured 3D morphing try to leverage pretrained 3D models to produce intermediate states. 
Interp3D~\cite{liu2026interp3d}, and MorphAny3D~\cite{sun2026morphany3d} build on TRELLIS~\cite{xiang2025trellis}, a pretrained 3D generation method, combining information from two source objects through predefined fusion strategies. Both these methods are training-free and do not learn a reusable morphing function optimized across a collection of example trajectories. However, generated sequences can exhibit artifacts, abrupt changes, or premature convergence toward one object. Moreover, representing source assets through 2D views can discard information about unobserved geometry
before morphing begins. %As a result, the quality of the generated transitions can depend on the viewpoint and quality of the input images.

Despite these limitations, training-free methods can provide a large number of pseudo-trajectories that would otherwise be unavailable for supervised learning. Our key idea is to use these trajectories as pseudo-supervision for learning a morphing function, rather than treating them as final outputs. By training across many source pairs and interpolation coefficients, we obtain a model that learns shared patterns of geometric and appearance transformations instead of relying on predefined fusion procedures. We hypothesize that learning these common patterns can improve generalization to unseen object pairs while reducing some of the irregularities present in the generated trajectories. %teacher-generated trajectories.

To realize this idea, we introduce \methodname{}, a teacher-student framework for textured 3D morphing. First, we use Interp3D as a teacher to generate morphing pseudo-targets of two source objects.
Then, we define a student model that adapts the TRELLIS generation process and learns to combine information from both sources using the teacher supervision. 
Unlike image-driven approaches, our framework conditions directly on the 3D representations of the sources, avoiding the need to reconstruct their geometry from images and the associated ambiguities in unobserved regions. 

However, teacher-generated pseudo-targets may still contain artifacts that could be transferred to the student during training. To mitigate this issue, we introduce a frozen TRELLIS-based generative prior that provides an additional regularization signal, encouraging the student to produce plausible intermediate states rather than merely reproducing the teacher's outputs.

% Our experiments show that \methodname{} produces intermediate morphing states with fewer artifacts and better preservation of source geometry w.r.t. previous training-free approaches. Additionally, the generated trajectories exhibit smoother and more balanced transitions, reducing abrupt changes and premature convergence toward either source object. To evaluate the progression of these transitions, we introduce Source Profile Error (SPE), a source-normalized metric that measures how perceptual changes are distributed along the morphing trajectory. Quantitatively, our method reduces PDV by 39\% compared to the strongest baseline, achieves the lowest SPE under multiple evaluation protocols, and obtains the best FID. Furthermore, once trained, our model requires no pair-specific optimization and generates morphing sequences approximately three times faster than existing approaches.
Evaluating morphing requires assessing both the quality of intermediate states and how the transition progresses between the sources. Measures of individual states or changes between consecutive states do not explicitly capture this source-relative progression. We therefore introduce Source Profile Error (SPE), a source-normalized metric that evaluates progression relative to both input objects. Our experiments against Interp3D and MorphAny3D show more regular perceptual progression and competitive visual quality, together with substantially lower generation time.

The main contributions of this paper are as follows:
\begin{itemize}
    \item We introduce \methodname{}, a teacher-student framework that learns textured 3D morphing from training-free pseudo-trajectories, without requiring exact ground-truth morphing sequences. \methodname{} is conditioned directly on the 3D representations of both source objects using a learned attention fusion mechanism called Gate-Adapted Morphing Cross-Attention (GAMCA).
    \item We introduce a generative regularization strategy based on a frozen TRELLIS prior, complementing teacher supervision to encourage plausible intermediate states.

    \item We propose Source Profile Error (SPE), a source-normalized metric that evaluates progression relative to both source assets, complementing existing quality and smoothness metrics.
\end{itemize}